% ##################################################################################################################
\section{Padang}
\label{sec:padang}
\hfill \textbf{Author:} Gregor Lämmel

% ##################################################################################################################
The Padang scenario demonstrates the application of \gls{matsim} to large-scale evacuation problems. The scenario was created as part of the third party funded project "Last-Mile". \citet{TaubenboeckEtAl2012ConcludingLastMilePaperNatHazards} gives a comprehensive overview.
Padang is locate at the west coast of Sumatra Island, Indonesia. As of 2014, the city has a population of about 1\,000\,000 people. 
Because if its prominent location at the coast in a so-called seismically locked area~\citep{McCloskey2010Padang2009Earthquake}, Padang is prone to earthquakes and subsequent tsunamis. Within the "Last-Mile" project a likely tsunami scenario triggered by an earthquake about 300\,km off the coast has been identified~\citep{GosebergSchlurmann2009HazardMappingPadang}. The assumed tsunami would leave about 30\,min for the evacuation. The flooding would reach as far as 3\,kilometers inland, and thus, it would endanger the life of up to 330\,000 people. \citet{Laemmel_PhDThesis_2011} developed a \gls{matsim} scenario representing the city with its affected population. The specialty of the situation in Padang is that it is expected everyone evacuates by foot. Simulating pedestrians with \gls{matsim} was a novelty at the time of the project. Therefore, the standard simulation model (see, e.g.,\,Chapter~\ref{sec:trafficflowmodel}) has been adapted in order to deal with pedestrians. 
Details are discussed in \citet{LaemmelKluepfelNagel2009EvacPadangAtBookTimmermanns}. Another important difference to most standard transport scenarios is that links in the network become flooded once the tsunami reaches them. Thus, the accessibility---flooded or not flood---for the links in the network is time-dependent. This time dependency is modeled by the concept of time-dependent networks \citep{LaemmelEtAl_TransResC_2010}. In the time-dependent network concept link attributes like \emph{freespeed} can by changed while the simulation is running by precomputed network change events. For the Padang scenario, the network change events have been extracted from microscopic flooding simulation data.
% ##################################################################################################################
%
\createfigure%
{Evacuation time analysis for downtown Padang. Numbers showing average evacuation time in minutes, which are also indicated by the colors green, yellow, red.}%
{Evacuation time analysis for downtown Padang. Numbers showing average evacuation time in minutes, which are also indicated by the colors green, yellow, red.}%
{\label{chap:using:padang}}%
{\includegraphics[width=0.6\textwidth, angle=0]{using/figures/dwntwnpdg}}%
{}

Key facts about the Padang scenario:
\begin{itemize}
\item The network consist of about 6\,000 nodes and 17\,000 links.
\item Synthetic populations for morning, afternoon, and night time have been created comprising up to 330\,000 agents.
\item The flooding is modeled by a set of 109\,network change events (one per minute) effecting 7\,609 links.
\item A set of 42\,shelter buildings, which can be used for vertical evacuation is also part of the scenario.
\end{itemize}
Based on the Padang scenario various evacuation strategies have been investigated:
\begin{itemize}
\item A seemingly obvious evacuation strategy is the shortest path solution, where everyone is on the shortest path. This solution, however, ignores possible congestion and lead to unfeasible results;
\item  An improvement is the Nash equilibrium approach, where everyone tries to find an optimal evacuation route by iterative learning \citep{LaemmelKluepfelNagel2009EvacPadangAtBookTimmermanns};
\item While the Nash equilibrium reduces the individual evacuation time, the total evacuation time might not be minimal. The marginal social cost based simulation approach tries to minimize the total evacuation time \citep{LaemmelFloetteroed_MertschingEtAl_2009,DresslerFloetteroedLaemmelNagelSkutella2010OptimalEvacuationLargeScaleScenarios};
\item This three basic evacuation approaches are investigated in combination with flooding \citep{LaemmelEtAl_TransResC_2010,Laemmel_PhDThesis_2011};
\item On top of this, an evacuation strategy that reduces the risk of exposure is developed by \citep{LaemmelKluepfelNagel2010PEDRiskPrinted};
\item And finally, \citet{FloetteroedLaemmel2010ICECShelterEvac} propose a method for integrating shelter buildings, which constitute sinks with limited capacity, into the simulation.  
\end{itemize}

% ##################################################################################################################
